% ============================================================
% SIGA — Documentación Técnica — Preámbulo compartido
% ============================================================
% Motor de compilación: pdflatex (Overleaf: Menu > Compiler > pdfLaTeX).
% REGLA para todo capítulo/apéndice: usar SOLO lo que ya está
% definido acá (paquetes, colores, estilos TikZ, comandos, entornos).
% No agregar \usepackage nuevos dentro de un capítulo — si hace
% falta algo que no está, se agrega ACÁ para que todos lo compartan.
% ============================================================

\documentclass[11pt,a4paper,oneside]{report}

% ---------- Idioma y codificación ----------
\usepackage[utf8]{inputenc}
\usepackage[T1]{fontenc}

% ---------- Ajuste fino de líneas ----------
% Comprime/expande el espaciado de caracteres unos puntos porcentuales para
% que más líneas encajen sin overfull \hbox, sin tocar ningún contenido.
% Reduce la mayoría de los "Overfull hbox" del documento (identificadores de
% código largos en \code{} contra el margen); no elimina los casos extremos,
% para esos hace falta un \allowbreak manual en el punto/guión bajo.
\usepackage{microtype}
% Ultimo recurso de TeX para armar un renglon: cuando ni microtype ni los
% \allowbreak alcanzan (parrafos con identificadores de codigo largos), permite
% estirar el espaciado de esa linea puntual en vez de dejar que el texto se
% salga del margen. Solo actua en los renglones que de otro modo darian
% Overfull \hbox; el resto del documento no cambia.
\setlength{\emergencystretch}{3em}

\usepackage[spanish]{babel}
% El módulo spanish activa < y > como shorthand de <<...>> (comillas
% angulares). Esto rompe cualquier texto con genéricos de C# (Result<T>,
% List<T>, HasConversion<int>()) dentro de argumentos móviles (TikZ,
% \section, \caption) con errores "extra }" / "runaway argument". Este
% documento no usa comillas angulares en ningún lado, así que se
% desactiva por completo en vez de escapar cada ocurrencia a mano.
% IMPORTANTE: babel solo activa los shorthands del idioma principal recién
% en \begin{document} (vía \selectlanguage interno) — un \shorthandoff{<>}
% suelto acá, en el preámbulo, se ejecuta ANTES de esa activación y queda
% pisado. Por eso se difiere con \AtBeginDocument, para que corra después.
\AtBeginDocument{\shorthandoff{<>}}
\usepackage{csquotes}

% ---------- Página ----------
\usepackage[margin=2.5cm,headheight=14pt]{geometry}
\usepackage{fancyhdr}
\pagestyle{fancy}
\fancyhf{}
\fancyhead[L]{\nouppercase{\leftmark}}
\fancyhead[R]{\thepage}
\renewcommand{\headrulewidth}{0.4pt}
\fancypagestyle{plain}{%
  \fancyhf{}
  \fancyhead[R]{\thepage}
  \renewcommand{\headrulewidth}{0pt}
}

% ---------- Títulos de capítulo ----------
\usepackage{titlesec}
\titleformat{\chapter}[display]
  {\normalfont\huge\bfseries}{\chaptertitlename\ \thechapter}{16pt}{\Huge}
\titlespacing*{\chapter}{0pt}{0pt}{32pt}

% ---------- Colores (paleta institucional de SIGA-Web, ver design-system.md) ----------
\usepackage{xcolor}
\definecolor{sigaAzul}{HTML}{00288E}
\definecolor{sigaCyan}{HTML}{006780}
\definecolor{sigaVerde}{HTML}{1B7A3D}
\definecolor{sigaAmbar}{HTML}{9A6700}
\definecolor{sigaRojo}{HTML}{B3261E}
\definecolor{sigaGris}{HTML}{5F6368}

% ---------- Bibliografía ----------
\usepackage[backend=biber,style=apa]{biblatex}
\addbibresource{bibliografia.bib}

% ---------- Enlaces internos/externos ----------
\usepackage{hyperref}
\hypersetup{
  colorlinks=true,
  linkcolor=sigaAzul,
  citecolor=sigaVerde,
  urlcolor=sigaCyan,
  pdftitle={SIGA-Optica -- Sistema Integral de Gestion y Agendamiento para Opticas},
  pdfauthor={Matias Melgarejo, Matias Gaona},
  bookmarksnumbered=true,
  bookmarksopen=true,
}

% ---------- Tablas ----------
\usepackage{booktabs}
\usepackage{longtable}
\usepackage{array}
\usepackage{tabularx}
\newcolumntype{Y}{>{\raggedright\arraybackslash}X}
\renewcommand{\arraystretch}{1.2}

% ---------- Listados de código ----------
\usepackage{listings}
\lstdefinestyle{sigaCode}{
  basicstyle=\ttfamily\small,
  keywordstyle=\color{sigaAzul}\bfseries,
  commentstyle=\color{sigaGris}\itshape,
  stringstyle=\color{sigaVerde},
  numbers=left,
  numberstyle=\tiny\color{sigaGris},
  numbersep=8pt,
  frame=single,
  framesep=4pt,
  breaklines=true,
  showstringspaces=false,
  tabsize=2,
  backgroundcolor=\color{black!3},
  captionpos=b,
  literate=
    {á}{{\'a}}1 {é}{{\'e}}1 {í}{{\'i}}1 {ó}{{\'o}}1 {ú}{{\'u}}1
    {Á}{{\'A}}1 {É}{{\'E}}1 {Í}{{\'I}}1 {Ó}{{\'O}}1 {Ú}{{\'U}}1
    {ñ}{{\~n}}1 {Ñ}{{\~N}}1 {ü}{{\"u}}1 {Ü}{{\"U}}1
    {¿}{{?`}}1 {¡}{{!`}}1,
}
\lstset{style=sigaCode}
% Nota: literate cubre tildes/ñ dentro de lstlisting/\lstinline. Los
% caracteres de dibujo de caja Unicode (├ └ │ ─) NO están cubiertos y
% NO hay que usarlos — para árboles de carpetas usar ASCII (|--, `--).

% ---------- Diagramas (TikZ) ----------
\usepackage{tikz}
\usetikzlibrary{positioning,arrows.meta,shapes.multipart,calc,fit,backgrounds,decorations.pathreplacing}

% --- Estilo compartido: diagramas de ARQUITECTURA / COMPONENTES ---
\tikzset{
  capa/.style={rectangle, rounded corners, draw=sigaAzul, thick,
    fill=sigaAzul!8, minimum width=3.2cm, minimum height=1.1cm,
    align=center, font=\small},
  subcapa/.style={rectangle, rounded corners, draw=sigaAzul!70, thin,
    fill=white, minimum width=2.6cm, minimum height=0.8cm,
    align=center, font=\scriptsize},
  externo/.style={rectangle, draw=sigaGris, thick, fill=sigaGris!10,
    minimum width=2.6cm, minimum height=1cm, align=center, font=\small},
  flechaArq/.style={-{Latex[length=2.5mm]}, thick, sigaGris},
}

% --- Estilo compartido: diagramas ER (entidad-relación) ---
% Uso: \node[entidad] (id) {\textbf{\textcolor{sigaAzul}{Nombre}} \nodepart{two} campo1 \\ campo2 \\ ...};
\tikzset{
  entidad/.style={
    rectangle split, rectangle split parts=2, draw=sigaAzul, thick,
    rectangle split part align={center,left}, rounded corners=1pt,
    fill=white, font=\footnotesize, text width=3.6cm, align=left,
    inner sep=4pt,
  },
  relacion/.style={-, thick, sigaGris},
  relacionN/.style={-{Bar[]}, thick, sigaGris},
}
% Título en negrita del compartimento superior de una entidad:
% primer renglón del nodo -> usar \textbf{\textcolor{sigaAzul}{Nombre}}

% --- Estilo compartido: diagramas de SECUENCIA ---
\tikzset{
  actor/.style={rectangle, draw=sigaAzul, thick, fill=sigaAzul!8,
    minimum width=2.6cm, minimum height=0.8cm, align=center,
    font=\small\bfseries},
  lineavida/.style={dashed, sigaGris},
  mensaje/.style={-{Latex[length=2mm]}, thick},
  mensajeretorno/.style={-{Latex[length=2mm]}, thick, dashed},
}

% ---------- Cajas de estado / advertencia / ADR ----------
\usepackage[most]{tcolorbox}
\tcbuselibrary{skins,breakable}

% Uso cajaAdvertencia/cajaRiesgo (el título es OPCIONAL, va con []):
%   \begin{cajaAdvertencia}  texto...  \end{cajaAdvertencia}                (título default "Atención"/"Riesgo")
%   \begin{cajaAdvertencia}[Título custom]  texto...  \end{cajaAdvertencia}
\newtcolorbox{cajaAdvertencia}[1][Atención]{
  colback=sigaAmbar!8, colframe=sigaAmbar, boxrule=0.8pt, arc=2pt,
  breakable, fonttitle=\bfseries\color{white}, colbacktitle=sigaAmbar,
  title=#1,
}
\newtcolorbox{cajaRiesgo}[1][Riesgo]{
  colback=sigaRojo!6, colframe=sigaRojo, boxrule=0.8pt, arc=2pt,
  breakable, fonttitle=\bfseries\color{white}, colbacktitle=sigaRojo,
  title=#1,
}
% Uso cajaADR (el título es OBLIGATORIO, va con {}):
%   \begin{cajaADR}{ADR 0007 -- El cristal ya no es Producto con stock}  texto...  \end{cajaADR}
\newtcolorbox{cajaADR}[1]{
  colback=sigaAzul!4, colframe=sigaAzul, boxrule=0.8pt, arc=2pt,
  breakable, fonttitle=\bfseries\color{white}, colbacktitle=sigaAzul,
  title=#1,
}

% ---------- Marcadores de estado (reemplazan ✅/⚠️/🟡 de Markdown) ----------
\newcommand{\Implementado}{\textcolor{sigaVerde}{\textbf{[Implementado]}}}
\newcommand{\Planificado}{\textcolor{sigaGris}{\textbf{[Planificado]}}}
\newcommand{\Atencion}{\textcolor{sigaAmbar}{\textbf{[Atención]}}}
\newcommand{\RiesgoTag}{\textcolor{sigaRojo}{\textbf{[Riesgo]}}}

% ---------- Código inline (reemplazo de `backtick` de Markdown) ----------
% Uso: \code{ver_pacientes}  ->  escapa _ % & # automáticamente, no hace
% falta escapar nada a mano adentro de \code{...}.
\newcommand{\code}[1]{\texttt{\detokenize{#1}}}

% ---------- Tabla de endpoints del Apéndice de Referencia de API ----------
% Evita repetir el boilerplate de longtable una vez por controller.
% Uso: \begin{apiTable}{\texttt{AuthController}}{tab:api-auth} ... \end{apiTable}
% OJO: #1 termina en un \caption{} (argumento móvil) -> usar \texttt{}, nunca \code{}.
\newenvironment{apiTable}[2]{%
  \setlength{\tabcolsep}{2pt}%
  \begin{longtable}{@{}p{1.3cm}p{5.0cm}p{3.2cm}p{2.9cm}p{2.9cm}@{}}
  \caption{#1}\label{#2}\\
  \toprule
  \textbf{Método} & \textbf{Ruta} & \textbf{Policy} & \textbf{Request} & \textbf{Response} \\
  \midrule
  \endfirsthead
  \toprule
  \textbf{Método} & \textbf{Ruta} & \textbf{Policy} & \textbf{Request} & \textbf{Response} \\
  \midrule
  \endhead
  \bottomrule
  \endfoot
}{%
  \end{longtable}
}

% ---------- Listas ----------
\usepackage{enumitem}
\setlist{itemsep=2pt, topsep=4pt}

% ---------- Figuras / captions ----------
\usepackage{graphicx}
\usepackage{caption}
\usepackage{subcaption}
\captionsetup{font=small, labelfont=bf, margin=1cm}

\linespread{1.08}

% ---------- Convención de \label (documentado, no forzado por código) ----------
% sec:<capitulo>-<seccion>   ej. sec:arch-multisucursal
% fig:<tipo>-<slug>          ej. fig:er-identidad, fig:seq-venta-a-pedido
% tab:<slug>                 ej. tab:api-auth, tab:enums
% mod:<numero>-<slug>        ej. mod:08-ventas
% adr:<numero>               ej. adr:0007
